%%% From the Foreign Trade Lecture

\section{International Law and Governance}

\begin{frame}{International Law}
    \mycomment{
    International law exists, but not in the same way as national, codified, positive law exists.

    There is natural law, which is considered to exist prior to set and codified law. Among this are the human rights, thought to exists independent of formulation. Set law that is agreed upon or issued by law makers is called positive law. 
    }
    \mytext{
        \begin{itemize}
        \item There is no World State
        \item International Law is common law / contract law
        \item Concept of Natural Law vs. positive law
        \item Declaration of Human Rights
        \item No enforcement authorities
        \item International Law of the Seas
        \item UN System
        \item OECD
        \item Local legal systems: EU Law
    \end{itemize}
    }
\end{frame}

\begin{frame}{International Law}
    \mycomment{
    International (positive) law originates from agreements between states. There is no super national authority that can force states under its will. All international cooperation is thus voluntary. However, there are standards that are agreed upon, constituting some kind of international common, or positive, law.
    }
    \mytext{
    \begin{itemize}
        \item Hugo Grotius 1583-1645
        \item States as subjects, not firms or people
        \item People to minor extent (per Art. 1 UN Charta)
        \item Also organisations of states
        \item UN is such an organisation
        \item UN charter
        \item UN declaration of human rights (non binding!)
        \item conventions and treaties (Geneva Convention Relating to the Status of Refugees)
        \item Specialised Agencies: IMF, World Bank, ILO etc.
    \end{itemize}    
    }
\end{frame}

\begin{frame}{International Law}
    \mycomment{
    In the economic realm, many frameworks exist to make international trade more safe and secure. 
    }
    \mytext{
     \begin{itemize}
        \item WTO $\neq$ UN, less members
        \item WTO $\leftarrow$ GATT
        \item GATT still in existence
        \item Dispute Settlement Body
        \item Not really active
        \item Art. 1 GATT: Most Favoured Nation Principle
        \item Exception: customs union 
    \end{itemize}   
    }
\end{frame}

\begin{frame}{WTO and Free Trade Areas}
    \showmytikz[1]{../shared/tikz/tikz-wto}
    \mycomment{
        The WTO is an organisation of member states based in Genevea. It runs, among others, a dispute settlement body. It came to replace the (still valid!) GATT, the general agreement on tariffs and trade.

        Within the WTO, the MFN most favoured nation principle holds: you are not allowed to charge anyone a tariff higher than you charge to the most favoured nation.

        Exceptions are only possible if you are within a free trade area. So no bilateral deals!
    }
\end{frame}

\begin{frame}{WTO and Free Trade Areas}
    \showmytikz[1,2]{../shared/tikz/tikz-wto}
    \mycomment{
        Since you can only positively discriminate your immediate neighbours if you are in a free trade area, free trade areas sprung up.

        The degree of economic integration within these areas varies widely. The simplest is just a free trade area, so no tariffs apply. But it can reach all the way to an integration like the EU, with a common currency, common standards, and common legal bodies.     

        (Picture is illustrative: USMCA has exactly three members, USA, Mexico, and Canada. Mercosur and EU have more members than shown, etc.)
    }
\end{frame}

\begin{frame}{WTO and Trade Blocks}
    \mytext{
    GATT/GATS/TRIPS $\rightarrow$ / WTO, Geneva: 
    \begin{itemize}
        \item `most favoured nation' clause
        \item ratchet principle (tarrifs can only go down)
        \item differentiation only justifiable by blocks
        \item result: surge in blocks
    \end{itemize}
    }
\end{frame}

\begin{frame}{Trade Blocks / PTAs }
    \showmytikz{../shared/tikz/tikz-ftas}
    \mycomment{
    Preferential trade agreements (PTAs) can be more or less intense. In a free trade agreement, you still need to file customs documents. In a customs union, you may still encounter different regulations. In a common market, you may still encounter different economic policies. In an economic union, you may still have to deal with exchange rates.     
    }
\end{frame}

\begin{frame}{Free Trade Zones Examples}
    \mycomment{
    There are plenty of free trade areas and monetary cooperation areas. Sometimes they go by French or English abbreviations. Negotiations typically take years, and scope, contracts, names and memberships may change.
    }
    \mytext{
    \begin{itemize}
        \item EMU, EU, EEA, EUCU
        \item Mercosur, CAN
        \item AFTA (ASEAN)
        \item JEFTA (EU-Japan), CETA (EU-Canada)
        \item CEDEAO=ECOWAS
        \item TTIP (?)
    \end{itemize}
    }
\end{frame}

\begin{frame}{Currency Blocks}
    \mytext{
    \begin{itemize}
        \item EMU (Euro)
        \item UEMOA=WAEMU and CEMAC (tied to \euro )
    \end{itemize}

    Unilateral peg
    \begin{itemize}
        \item Kosovo
        \item Zimbabwe
        \item Morocco
        \item Argentina..?
    \end{itemize}
    }
\end{frame}


\begin{frame}{Regional, Economic and Political Cooperation}
    \mytext{
     \begin{itemize}
        \item Working Groups (G7, G20)
        \item Local Working Groups (E3)
        \item Currency union(s), customs union(s)
        \item EU, EEA, ECU
        \item MERCOSUR
        \item ASEAN
        \item USMCA
        \item Economic Groups 
        \item OPEC
        \item BRICS
        \item ...
    \end{itemize}   
    }
    \mycomment{
    The intensity of the cooperation can and does vary widely. Most of them are mere conferences, such as the G7, that may not be more than an office in New York at the UN. 

    In others, most prominently the European Union, member states parts of of their sovereingty to multi-national bodies with large institutions - and budgets.
    }
\end{frame}



\begin{frame}{Freedom of the Seas}
    \mytext{
    \begin{itemize}
        \item mare liberum (Hugo Grotius)
        \item 17th century
        \item British concept, 3 nautical miles (5.6 km)
        \item later de-facto 12 nautical miles
        \item ships under jurisdiction of flag they fly
    \end{itemize}
    }
\end{frame}

% \begin{frame}{The Gambia}
%     \mycomment{
%         To extend a rule from land to the sea just as far as you can look and fire a canon was sometimes even used the other way round. 
        
%         Look at the shape of The Gambia. It extends six miles left and right of the river Gambia---just as far as a British war ship was able to fire at the time. 
        
%         The Gambia is totally surrounded by Senegal, then under French rule. Still today, they use English and French as official languages, and have legal systems more leaning towards French written or English common law.
%     }
%     \showmybitmap[https://de.wikipedia.org/wiki/Datei:Gambia_rel_location_map.svg]{shared/bitmaps/Gambia_rel_location_map.svg.png}
% \end{frame}

\begin{frame}{UNCLOS}
    \mytext{
    \begin{itemize}
        \item Four treaties, 1958/1962
        \item 63 states, NATO, Soviet-bloc, OPEC etc.
        \item today 169 states, but not US
        \item Convention on the Territorial Sea and Continguous Zone
        \item Convention on the Continental Shelf
        \item Convention on the High Seas
        \item Convention on Fishing and Convention of Living Resources of the High Seas
        \item United Nations Convention on the Law of the Sea
    \end{itemize}
    }
\end{frame}

% \begin{frame}{International Waters and Exclusive Zones}
%     \showmytikz{../shared/tikz/tikz-waters.tex}
%     \mycomment{
%     There is a difference between exclusive right to economic use (of the waters, e.g. for fishing, or the resources at the sea bed) and legislation in regards to ships and traffic.

%     The legislation on board of the ship depends on  the flag under which the ship sails. However, contracts may specify which legislation applies, e.g. for work contracts.

%     \begin{equation}
%         1 NM = 1.852km
%     \end{equation}
%     A nautical mile is approximately one minute of arc of latitude.
%     }
% \end{frame}

\begin{frame}{International Waters and Exclusive Zones}
    \showmytikz{../shared/tikz/tikz-waters-two.tex}
    \mycomment{
    The territorial waters---and the contiguous zone---are ruled by the state just like its territory.

    Beyond that lies the freedom of the seas, and vessels have free passage. However, the state reserves the right to the sea bed, that is to lay cables, pipelines, drilling and so forth. This extends to the end of the continental shelf or to 200NM, whichever is greater.

    The exact shape and extend of both EEZ, the end of the continental shelf, and the extended continental shelf can give reason to disputes and conflict. Tiny islands often come with the benefit of large undersea territories.
    }
\end{frame}


\begin{frame}{UN System}
    \showmytikz{../shared/tikz/tikz-unsystem}
    \mycomment{
         The UN is an organisation with states as members.
        
        It has six principal organs:
        \begin{itemize}
            \item General Assembly
            \item Secretariat
            \item International Court of Justice
            \item Security Council
            \item Economic and Social Coucil
            \item (Trusteeship Council)
        \end{itemize}

    There are many specialised bodys and programmes:
    
    \begin{itemize}
        \item Funds and programmes (UNDP, UNICEF etc.)
        \item Research and training institutes (UNITAR, UNU etc.)
        \item Other entities and bodies (UNCTAD, UNHCR, UNIFEM etc.)
    \end{itemize}

    }
\end{frame}

\begin{frame}{EU System}
    \showmytikz{../shared/tikz/tikz-eusystem}
    \mycomment{
    The constitutional text are the Treaty of the European Union (TEU) and the Treaty of the Functioning of the European Union (TFEU).

    \begin{itemize}
        \item Regulations: directly binding
        \item Directives: must be implemented by member states
        \item Communications: non-binding
    \end{itemize}
    }
\end{frame}

\begin{frame}{European Union - Competencies}
    \showmytikz{../shared/tikz/tikz-eucompetencies}
    \mycomment{The competencies of the European Union are defined in Art. 2 TFEU, while Art. 3 lists the exclusive competencies explicitly.

    Competition rules are an exclusive competency as long as the single market is concerned. 
    
    Monetary policy is an exclusive competency in regards to the Eurozone countries. 
    
    Trade is the common commercial policy. There are exceptions, in particular in areas where all competencies remain with member states (e.g. military equipment). 
    
    Fishery is an exclusive competency, but only as far as it is regulated by the common fishery policy. 
    }
\end{frame}



\begin{frame}{OECD System}
    \showmytikz{../shared/tikz/tikz-oecdsystem}
    \mycomment{
     Organisation for Economic Cooperation and Development
    \begin{itemize}
        \item Not an organisation, but a permanent conference (Paris)
        \item Organs and specialised agencies:
        \item International Energy Agency IEA
    \end{itemize}   
    }
\end{frame}

\begin{frame}{Security and Military Alliances}
    \mytext{
    \begin{itemize}
        \item OSCE Organisation for Security and Co-operation in Europe (Vienna)
        \begin{itemize}
            \item All states in Europe  
            \item Conference on security
        \end{itemize}
        \item NATO 
        \item Council of Europe
            \begin{itemize}
                \item Not part of EU system
                \item European Court of Human Rights
                \item European Convention on Human Rights
            \end{itemize}        
    \end{itemize}
    }
\end{frame}



\begin{frame}{Asia/Pacific: ASEAN}
    \mycomment{
    Asian countries are typically dollarized, as they typically trade more with the US (on the other side of the Pacific Ocean) than with Europe. Economic integration in the area is comparatively weak. 
    
    Japan in particular is a single country with a single currency (Japanese Yen, which means "unit") without a "Yen zone". 
    
    The Chinese economy is huge, but capital flows are typically restricted; the currency is the "Renminbi", which means "currency of the people", with the unit Yüan, which also means "unit".
    }
    \mytext{
    ASEAN Association of Southeat Asian Nations
    \begin{itemize}
        \item AFTA Free trade area
        \item AIA Asean Investment Area
    \end{itemize}

    }
\end{frame}

\begin{frame}{Africa: AEC, regional blocks}
    \mycomment{
    Africa with 56 states has a plethora of international organisations. Francophone Africa is closely tied to the Euro area, while chiefly anglophone South Africa (the Rand zone) and East Africa are more leaning towards the US Dollar.
    }
    \mytext{
        \begin{itemize}
        \item AEC African Economic Community
        \begin{itemize}
            \item ECOWAS/CEDEAO Economic Community of West African States
            \item WAEMU/UEMOA West African Economic and Monetary Union (Franc CFA)
            \item ECCAS/CEEAC Economic Community of Central African States    
            \item EMCCA/CEMAC Economic and Monetary Community of Central Africa
            \item ...
            \item SADC Southern African Development Community
            \item SACU Southern African Customs Union
            \item ...
        \end{itemize}
    \end{itemize}
    }
\end{frame}

